\begingroup\setbox0=\vbox{\section*{Triangles and Their Area}}\endgroup
\chapterhead{11}{Triangles and Their Area}

\shead{Base, height and area}
\begin{knowbox}
\begin{minipage}[t]{0.6\linewidth}
A \textbf{triangle} has $3$ sides, $3$ vertices and $3$ angles. By its angles it is \emph{acute} (three acute angles), \emph{obtuse} (one obtuse angle) or \emph{right-angled} (one right angle).\par
\smallskip
The segment drawn from a vertex \textbf{perpendicular} to the opposite side is a \textbf{height}; that side is the \textbf{base}. Every triangle has $3$ heights.\par
\smallskip
The \textbf{area of a triangle} is the base times the height (in the same unit), divided by $2$:
\[S=\frac{a\times h}{2}\qquad h=\frac{S\times2}{a}\qquad a=\frac{S\times2}{h}\]
In a right-angled triangle, the two sides at the right angle are each other's base and height.
\end{minipage}\hfill
\begin{minipage}[t]{0.36\linewidth}\vspace{2pt}
\centering
\begin{tikzpicture}[scale=0.9]
\coordinate (B) at (0,0); \coordinate (C) at (4,0); \coordinate (A) at (1.3,2.6); \coordinate (H) at (1.3,0);
\fill[tmBlue!12] (A)--(B)--(C)--cycle;
\draw[geo] (A)--(B)--(C)--cycle;
\draw[dashed,line width=0.7pt,tmOrange!80!black] (A)--(H);
\ramark{H}{C}{A}
\node[lbl,above] at (A){A}; \node[lbl,below left] at (B){B}; \node[lbl,below right] at (C){C}; \node[lbl,below] at (H){H};
\node[dim,right] at (1.3,1.3){$h$};
\draw[|<->|,Slate] (0,-0.55)--node[fill=white,dim]{$a$}(4,-0.55);
\end{tikzpicture}
\end{minipage}
\end{knowbox}
\gap
\begin{exbox}{\ExWord\ 1}
Find the area of a triangle with:
\begin{cols2}
\pa{a} base $8$ cm, height $5$ cm; & \pa{b} base $2.5$ dm, height $12$ cm.
\end{cols2}
\guide
\pa{a} $S=8\times5:2=20$ cm$^2$.\par
\pa{b} Use the same unit first: $2.5$ dm $=25$ cm.\par $S=25\times12:2=150$ cm$^2$.
\dapso{\pa{a} $20$ cm$^2$;\par \pa{b} $150$ cm$^2$.}
\end{exbox}
\gap
\begin{exbox}{\ExWord\ 2}
\begin{minipage}[t]{0.62\linewidth}
Triangle ABC has a right angle at A, with $AB=6$ cm, $AC=8$ cm and $BC=10$ cm. Find its area and the length of the height AH.
\end{minipage}\hfill
\begin{minipage}[t]{0.34\linewidth}\vspace{0pt}\centering
\begin{tikzpicture}[scale=0.36]
\coordinate (A) at (0,0); \coordinate (B) at (0,6); \coordinate (C) at (8,0); \coordinate (H) at (2.88,3.84);
\draw[geo] (A)--(B)--(C)--cycle; \ramark{A}{B}{C}
\draw[dashed,line width=0.7pt,tmOrange!80!black] (A)--(H); \ramark{H}{C}{A}
\node[lbl,below left] at (A){A}; \node[lbl,above] at (B){B}; \node[lbl,below right] at (C){C}; \node[lbl,above right] at (H){H};
\end{tikzpicture}
\end{minipage}
\guide
Take AC as the base; then AB is the height: $S=6\times8:2=24$ cm$^2$.\par
Take BC as the base; then AH is the height: $AH=24\times2:10=4.8$ cm.
\dapso{$24$ cm$^2$;\par $AH=4.8$ cm.}
\end{exbox}
\gap
\begin{exbox}{\ExWord\ 3}
A triangle has area $36$ cm$^2$ and base $9$ cm. Find its height.
\guide
$h=S\times2:a=36\times2:9=8$ cm.
\dapso{$8$ cm.}
\end{exbox}
\gap
\begin{exbox}{\ExWord\ 4}
\begin{minipage}[t]{0.6\linewidth}
Rectangle ABCD has $AB=12$ cm and $BC=8$ cm; M is the midpoint of BC. Find the areas of triangles ABM, MCD and AMD.
\end{minipage}\hfill
\begin{minipage}[t]{0.36\linewidth}\vspace{0pt}\centering
\begin{tikzpicture}[scale=0.3]
\coordinate (A) at (0,8); \coordinate (B) at (12,8); \coordinate (C) at (12,0); \coordinate (D) at (0,0); \coordinate (M) at (12,4);
\fill[shade] (A)--(M)--(D)--cycle;
\draw[geo] (A)--(B)--(C)--(D)--cycle; \draw[geo] (A)--(M)--(D);
\node[lbl,above left] at (A){A}; \node[lbl,above right] at (B){B}; \node[lbl,below right] at (C){C}; \node[lbl,below left] at (D){D}; \node[lbl,right] at (M){M};
\end{tikzpicture}
\end{minipage}
\guide
$BM=MC=8:2=4$ cm.\par $S_{ABM}=12\times4:2=24$ cm$^2$;\par $S_{MCD}=12\times4:2=24$ cm$^2$.\par
$S_{ABCD}=12\times8=96$ cm$^2$, so $S_{AMD}=96-24-24=48$ cm$^2$.\par
\emph{Remark:} triangle AMD has base AD and height equal to AB, so its area is half the rectangle.
\dapso{$24$ cm$^2$; $24$ cm$^2$; $48$ cm$^2$.}
\end{exbox}
\dividerbar{Practice}
\begin{probbox}{1}{\relax}
Find the area of a triangle with base $a$ and height $h$:
\begin{cols2}
\pa{a} $a=15$ cm, $h=8$ cm; & \pa{b} $a=4.5$ m, $h=2$ m;\\
\pa{c} $a=\dfrac34$ m, $h=\dfrac25$ m; & \pa{d} $a=1.2$ dm, $h=15$ cm.
\end{cols2}
\end{probbox}
\probsol{
\pa{a} $15\times8:2=60$ cm$^2$.\par \pa{b} $4.5\times2:2=4.5$ m$^2$.\par \pa{c} $\dfrac34\times\dfrac25:2=\dfrac3{20}$ m$^2$.\par \pa{d} $1.2$ dm $=12$ cm; $12\times15:2=90$ cm$^2$.
\dapso{$60$ cm$^2$; $4.5$ m$^2$; $\dfrac3{20}$ m$^2$; $90$ cm$^2$.}}
\gap
\begin{probbox}{2}{\relax}
Find the area of a right-angled triangle whose two shorter sides are $3.5$ dm and $2$ dm.
\end{probbox}
\probsol{
$S=3.5\times2:2=3.5$ dm$^2$.
\dapso{$3.5$ dm$^2$.}}
\gap
\begin{probbox}{3}{\relax}
A triangle has area $45$ cm$^2$ and height $6$ cm. Find its base.
\end{probbox}
\probsol{
$a=45\times2:6=15$ cm.
\dapso{$15$ cm.}}
\gap
\begin{probbox}{4}{\relax}
A triangular plot has base $40$ m and height $\dfrac34$ of the base. Find its area.
\end{probbox}
\probsol{
Height: $40\times\dfrac34=30$ m.\par Area: $40\times30:2=600$ m$^2$.
\dapso{$600$ m$^2$.}}
\gap
\begin{probbox}{5}{\relax}
A triangular card has base $2.4$ dm and height $\dfrac23$ of the base. Find its area.
\end{probbox}
\probsol{
Height: $2.4\times\dfrac23=1.6$ dm.\par Area: $2.4\times1.6:2=1.92$ dm$^2$.
\dapso{$1.92$ dm$^2$.}}
\gap
\begin{probbox}{6}{\relax}
Triangle ABC has a right angle at A, $AB=9$ cm, $AC=12$ cm and $BC=15$ cm. Find the height AH.
\end{probbox}
\probsol{
$S=9\times12:2=54$ cm$^2$.\par $AH=54\times2:15=7.2$ cm.
\dapso{$7.2$ cm.}}
\gap
\begin{probbox}{7}{\relax}
A triangle has the same area as a square of side $6$ cm. Its base is $9$ cm. Find its height.
\end{probbox}
\probsol{
Area: $6\times6=36$ cm$^2$.\par Height: $36\times2:9=8$ cm.
\dapso{$8$ cm.}}
\gap
\begin{probbox}{8}{\relax}
A triangle has base $18$ cm. Lengthening the base by $4$ cm (same height) adds $20$ cm$^2$ to the area. Find the original area.
\end{probbox}
\probsol{
The extra part is a triangle with base $4$ cm and area $20$ cm$^2$; its height is $20\times2:4=10$ cm.\par
Original area: $18\times10:2=90$ cm$^2$.
\dapso{$90$ cm$^2$.}}
\gap
\begin{probbox}{9}{\relax}
Square ABCD has side $10$ cm; M is the midpoint of BC. Find the areas of triangle ABM and quadrilateral AMCD.
\end{probbox}
\probsol{
$S_{ABM}=10\times5:2=25$ cm$^2$.\par $S_{AMCD}=10\times10-25=75$ cm$^2$.
\dapso{$25$ cm$^2$; $75$ cm$^2$.}}
\gap
\begin{probbox}{10}{\relax}
Triangle ABC has area $60$ cm$^2$ and $BC=15$ cm. Point M lies on BC with $BM=5$ cm. Find the area of triangle ABM.
\end{probbox}
\probsol{
The height from A to BC is $60\times2:15=8$ cm.\par Triangle ABM has the same height: $S_{ABM}=5\times8:2=20$ cm$^2$.
\dapso{$20$ cm$^2$.}}
\gap
\begin{probbox}{11}{\relax}
\begin{minipage}[t]{0.62\linewidth}
How many triangles are there in the figure? Name them.
\end{minipage}\hfill
\begin{minipage}[t]{0.34\linewidth}\vspace{0pt}\centering
\begin{tikzpicture}[scale=0.6]
\coordinate (A) at (2,2.6); \coordinate (B) at (0,0); \coordinate (M) at (1.3,0); \coordinate (N) at (2.6,0); \coordinate (C) at (4,0);
\draw[geo] (A)--(B)--(C)--cycle; \draw[geo] (A)--(M) (A)--(N);
\node[lbl,above] at (A){A}; \node[lbl,below] at (B){B}; \node[lbl,below] at (M){M}; \node[lbl,below] at (N){N}; \node[lbl,below] at (C){C};
\end{tikzpicture}
\end{minipage}
\end{probbox}
\probsol{
Each triangle has vertex A and a base joining two of the four points B, M, N, C.\par There are $3+2+1=6$ such segments.\par
Triangles: ABM, AMN, ANC, ABN, AMC, ABC.
\dapso{$6$ triangles.}}
\gap
\begin{probbox}{12}{\relax}
Find the area of a triangle with base $14$ cm and height $9$ cm.
\end{probbox}
\probsol{
$S=14\times9:2=63$ cm$^2$.
\dapso{$63$ cm$^2$.}}
\gap
\begin{probbox}{13}{\relax}
Find the area of a triangle with base $2.5$ m and height $1.6$ m.
\end{probbox}
\probsol{
$S=2.5\times1.6:2=2$ m$^2$.
\dapso{$2$ m$^2$.}}
\gap
\begin{probbox}{14}{\relax}
Find the area of a triangle with base $\dfrac34$ m and height $\dfrac25$ m.
\end{probbox}
\probsol{
$S=\dfrac34\times\dfrac25:2=\dfrac3{20}$ m$^2$.
\dapso{$\dfrac3{20}$ m$^2$.}}
\gap
\begin{probbox}{15}{\relax}
Find the area of a triangle with base $3$ dm and height $25$ cm.
\end{probbox}
\probsol{
$3$ dm $=30$ cm.\par $S=30\times25:2=375$ cm$^2$.
\dapso{$375$ cm$^2$.}}
\gap
\begin{probbox}{16}{\relax}
Find the area of a right triangle whose legs are $7$ cm and $10$ cm.
\end{probbox}
\probsol{
One leg is the base and the other the height.\par $S=7\times10:2=35$ cm$^2$.
\dapso{$35$ cm$^2$.}}
\gap
\begin{probbox}{17}{\relax}
A triangle has area $54$ cm$^2$ and base $12$ cm. Find its height.
\end{probbox}
\probsol{
Height $=$ area $\times2\,:$ base.\par $h=54\times2:12=9$ cm.
\dapso{$9$ cm.}}
\gap
\begin{probbox}{18}{\relax}
A triangle has area $7.5$ dm$^2$ and height $2.5$ dm. Find its base.
\end{probbox}
\probsol{
Base $=$ area $\times2\,:$ height.\par $a=7.5\times2:2.5=6$ dm.
\dapso{$6$ dm.}}
\gap
\begin{probbox}{19}{\relax}
A triangular plot has base $50$ m and height $24$ m. Each square metre yields $0.8$ kg of vegetables. What is the total yield?
\end{probbox}
\probsol{
Area: $50\times24:2=600$ m$^2$.\par Yield: $0.8\times600=480$ kg.
\dapso{$480$ kg.}}
\gap
\begin{probbox}{20}{\relax}
A triangle has base $16$ cm and height $\dfrac34$ of the base. Find its area.
\end{probbox}
\probsol{
Height: $16\times\dfrac34=12$ cm.\par Area: $16\times12:2=96$ cm$^2$.
\dapso{$96$ cm$^2$.}}
\gap
\begin{probbox}{21}{\relax}
The base and height of a triangle add up to $30$ cm, and the base is $6$ cm longer. Find the area.
\end{probbox}
\probsol{
Base: $(30+6):2=18$ cm.\par Height: $12$ cm.\par Area: $18\times12:2=108$ cm$^2$.
\dapso{$108$ cm$^2$.}}
\gap
\begin{probbox}{22}{\relax}
Triangle ABC has a right angle at A, $AB=15$ cm, $AC=20$ cm, $BC=25$ cm. Find its area and the height AH.
\end{probbox}
\probsol{
Area: $15\times20:2=150$ cm$^2$.\par With BC as base, AH is the height.\par $AH=150\times2:25=12$ cm.
\dapso{$150$ cm$^2$; $AH=12$ cm.}}
\gap
\begin{probbox}{23}{\relax}
Triangle one has base $12$ cm and height $8$ cm. Triangle two has the same area and base $16$ cm. Find its height.
\end{probbox}
\probsol{
Area: $12\times8:2=48$ cm$^2$.\par Height of triangle two: $48\times2:16=6$ cm.
\dapso{$6$ cm.}}
\gap
\begin{probbox}{24}{\relax}
Each triangular flag has base $30$ cm and height $20$ cm. How many square decimetres of cloth do $20$ flags need? (Ignore seams.)
\end{probbox}
\probsol{
One flag: $30\times20:2=300$ cm$^2$.\par $20$ flags: $6\,000$ cm$^2$ $=60$ dm$^2$.
\dapso{$60$ dm$^2$.}}
\gap
\begin{probbox}{25}{\relax}
A right triangle has area $24$ cm$^2$ and one leg $6$ cm. Find the other leg.
\end{probbox}
\probsol{
The other leg is the height for the base $6$ cm.\par Length: $24\times2:6=8$ cm.
\dapso{$8$ cm.}}
\gap
\begin{probbox}{26}{\relax}
Rectangle ABCD is $12$ cm by $8$ cm. Diagonal AC cuts it into two triangles. Find the area of each.
\end{probbox}
\probsol{
Triangle ABC has a right angle at B with legs $12$ cm and $8$ cm.\par $S_{ABC}=12\times8:2=48$ cm$^2$.\par Likewise $S_{ACD}=48$ cm$^2$ (half the rectangle).
\dapso{$48$ cm$^2$ each.}}
\gap
\begin{probbox}{27}{\relax}
A triangle has area $150$ cm$^2$. Raising its height by $3$ cm (same base) adds $30$ cm$^2$. Find the base and the height.
\end{probbox}
\probsol{
The extra part is a triangle with the same base, height $3$ cm and area $30$ cm$^2$.\par Base: $30\times2:3=20$ cm.\par Height: $150\times2:20=15$ cm.
\dapso{Base $20$ cm, height $15$ cm.}}
\gap
\begin{probbox}{28}{\relax}
A triangle with base $8$ cm has the same area as a square of side $8$ cm. Find its height.
\end{probbox}
\probsol{
Square: $8\times8=64$ cm$^2$.\par Height: $64\times2:8=16$ cm.
\dapso{$16$ cm.}}
\gap
\begin{probbox}{29}{\relax}
A triangular field has base $120$ m and height $80$ m. Give its area in hectares.
\end{probbox}
\probsol{
Area: $120\times80:2=4\,800$ m$^2$.\par $1$ ha $=10\,000$ m$^2$, so this is $0.48$ ha.
\dapso{$0.48$ ha.}}
\gap
\begin{probbox}{30}{\relax}
Triangle ABC has area $48$ cm$^2$; M is the midpoint of BC. Find the area of triangle ABM.
\end{probbox}
\probsol{
ABM and ABC share the height from A.\par $BM=\dfrac12BC$, so $S_{ABM}=\dfrac12S_{ABC}$.\par $S_{ABM}=48:2=24$ cm$^2$.
\dapso{$24$ cm$^2$.}}
\dividerbar{Challenge -- Extension}
\begin{probbox}{31}{\stars{1}}
A triangle has area $120$ cm$^2$. Shortening the base by $3$ cm (same height) takes $18$ cm$^2$ off the area. Find the base.
\end{probbox}
\probsol{
Height: $18\times2:3=12$ cm.\par Base: $120\times2:12=20$ cm.
\dapso{$20$ cm.}}
\gap
\begin{probbox}{32}{\stars{1}}
The two shorter sides of a right-angled triangle add up to $14$ cm, and one is $2$ cm longer than the other. Find its area.
\end{probbox}
\probsol{
The sides are $(14+2):2=8$ cm and $6$ cm.\par Area: $8\times6:2=24$ cm$^2$.
\dapso{$24$ cm$^2$.}}
\gap
\begin{probbox}{33}{\stars{1}}
Rectangle ABCD has $AB=16$ cm and $AD=10$ cm. Point E is on AB with $AE=6$ cm. Find the areas of triangles ADE, EBC and ECD.
\end{probbox}
\probsol{
$S_{ADE}=6\times10:2=30$ cm$^2$;\par $EB=10$ cm, so $S_{EBC}=10\times10:2=50$ cm$^2$.\par
$S_{ECD}=16\times10:2=80$ cm$^2$ (half the rectangle).\par \emph{Check:} $30+50+80=160=16\times10$.
\dapso{$30$ cm$^2$; $50$ cm$^2$; $80$ cm$^2$.}}
\gap
\begin{probbox}{34}{\stars{2}}
Triangle ABC has a right angle at A, $AB=6$ cm and $AC=8$ cm. M is on AB with $AM=2$ cm, and N is on AC with $AN=4$ cm. Find the area of quadrilateral BMNC.
\end{probbox}
\probsol{
$S_{ABC}=6\times8:2=24$ cm$^2$;\par $S_{AMN}=2\times4:2=4$ cm$^2$.\par $S_{BMNC}=24-4=20$ cm$^2$.
\dapso{$20$ cm$^2$.}}
\gap
\begin{probbox}{35}{\stars{2}}
\begin{minipage}[t]{0.62\linewidth}
How many triangles are there in the figure?
\end{minipage}\hfill
\begin{minipage}[t]{0.34\linewidth}\vspace{0pt}\centering
\begin{tikzpicture}[scale=0.6]
\coordinate (A) at (2,3); \coordinate (B) at (0,0); \coordinate (C) at (4,0);
\draw[geo] (A)--(B)--(C)--cycle;
\foreach \x in {1,2,3}{\draw[geo] (A)--(\x,0);}
\draw[geo] (1,1.5)--(3,1.5);
\end{tikzpicture}
\end{minipage}
\end{probbox}
\probsol{
Ignoring the middle crossbar: the base has $5$ points, giving $4+3+2+1=10$ triangles with vertex A.\par
The crossbar also has $5$ points, giving $10$ more triangles with vertex A.\par In all: $20$ triangles.
\dapso{$20$ triangles.}}
\gap
\begin{probbox}{36}{\stars{1}}
In triangle ABC, $BC=20$ cm. M lies on BC with $BM=8$ cm, and triangle ABM has area $40$ cm$^2$. Find the area of ABC.
\end{probbox}
\probsol{
The common height from A: $40\times2:8=10$ cm.\par $S_{ABC}=20\times10:2=100$ cm$^2$.
\dapso{$100$ cm$^2$.}}
\gap
\begin{probbox}{37}{\stars{1}}
Rectangle ABCD has $AB=20$ cm, $AD=12$ cm. M is the midpoint of AB and N the midpoint of BC. Find the area of triangle DMN.
\end{probbox}
\probsol{
Rectangle: $20\times12=240$ cm$^2$.\par $S_{AMD}=60$, $S_{MBN}=30$, $S_{NCD}=60$ cm$^2$.\par $S_{DMN}=240-60-30-60=90$ cm$^2$.
\dapso{$90$ cm$^2$.}}
\gap
\begin{probbox}{38}{\stars{1}}
Square ABCD has side $12$ cm, and E is any point on AB. Find the area of triangle ECD and the total area of triangles AED and EBC.
\end{probbox}
\probsol{
Triangle ECD has base $CD=12$ cm and height $12$ cm.\par $S_{ECD}=72$ cm$^2$.\par The other two total $144-72=72$ cm$^2$, wherever E is.
\dapso{Both are $72$ cm$^2$.}}
\gap
\begin{probbox}{39}{\stars{1}}
A triangle has area $40$ cm$^2$. If its base is tripled and its height halved, what is the new area?
\end{probbox}
\probsol{
Tripling the base triples the area.\par Halving the height halves it.\par New area: $40\times3:2=60$ cm$^2$.
\dapso{$60$ cm$^2$.}}
\gap
\begin{probbox}{40}{\stars{1}}
The legs of a right triangle add up to $23$ cm and differ by $7$ cm. Find its area.
\end{probbox}
\probsol{
Longer leg: $(23+7):2=15$ cm.\par Shorter: $8$ cm.\par Area: $15\times8:2=60$ cm$^2$.
\dapso{$60$ cm$^2$.}}
\gap
\begin{probbox}{41}{\stars{1}}
Triangle ABC has area $90$ cm$^2$. M on BC makes $MC$ twice $BM$. Find the areas of ABM and AMC.
\end{probbox}
\probsol{
$BM=\dfrac13BC$ and $MC=\dfrac23BC$; both share the height from A.\par $S_{ABM}=90:3=30$ cm$^2$.\par $S_{AMC}=60$ cm$^2$.
\dapso{$S_{ABM}=30$ cm$^2$, $S_{AMC}=60$ cm$^2$.}}
\gap
\begin{probbox}{42}{\stars{1}}
Triangle ABC has a right angle at A, $AB=12$ cm, $AC=16$ cm. M is the midpoint of AC. Compare the areas of ABM and BMC.
\end{probbox}
\probsol{
$AM=MC=8$ cm; both have height $AB=12$ cm.\par $S_{ABM}=48$ cm$^2$.\par $S_{BMC}=48$ cm$^2$.\par They are equal.
\dapso{Equal, $48$ cm$^2$ each.}}
\gap
\begin{probbox}{43}{\stars{1}}
Triangle ABC has $BC=24$ cm and height $AH=10$ cm. BC is extended beyond C by $CD=6$ cm. Find the areas of ACD and ABD.
\end{probbox}
\probsol{
ABC, ACD and ABD all have height $AH=10$ cm.\par $S_{ACD}=6\times10:2=30$ cm$^2$.\par $S_{ABD}=30\times10:2=150$ cm$^2$.
\dapso{$30$ cm$^2$ and $150$ cm$^2$.}}
\gap
\begin{probbox}{44}{\stars{1}}
Two triangles have equal heights. The base of the first is $\dfrac23$ of the second's, and the second's area is $24$ cm$^2$ larger. Find each area.
\end{probbox}
\probsol{
With equal heights, areas are in the ratio of the bases: $S_1=\dfrac23S_2$.\par The difference of $3-2=1$ part is $24$ cm$^2$.\par $S_1=48$ cm$^2$; $S_2=72$ cm$^2$.
\dapso{$48$ cm$^2$ and $72$ cm$^2$.}}
\gap
\begin{probbox}{45}{\stars{1}}
A garden is $40$ m by $25$ m. Flowers grow on a triangle whose base is one long side and whose apex is on the opposite side; grass covers the rest. Find both areas.
\end{probbox}
\probsol{
The triangle has base $40$ m and height $25$ m.\par Flowers: $500$ m$^2$.\par Garden: $1\,000$ m$^2$.\par Grass: $500$ m$^2$.
\dapso{Flowers $500$ m$^2$, grass $500$ m$^2$.}}
\gap
\begin{probbox}{46}{\stars{1}}
Triangle ABC has area $72$ cm$^2$. D is on AB with $AD=\dfrac13AB$; E is the midpoint of AC. Find the area of ADE.
\end{probbox}
\probsol{
Join D to C.\par ADC and ABC share the height from C, and $AD=\dfrac13AB$.\par $S_{ADC}=72:3=24$ cm$^2$.\par ADE and ADC share the height from D, and $AE=\dfrac12AC$.\par $S_{ADE}=12$ cm$^2$.
\dapso{$12$ cm$^2$.}}
\gap
\begin{probbox}{47}{\stars{1}}
Square ABCD has side $10$ cm. E is the midpoint of BC and F the midpoint of CD. Find the area of triangle AEF.
\end{probbox}
\probsol{
$S_{ABE}=25$ cm$^2$; $S_{ADF}=25$ cm$^2$.\par $S_{ECF}=5\times5:2=12.5$ cm$^2$.\par $S_{AEF}=100-25-25-12.5=37.5$ cm$^2$.
\dapso{$37.5$ cm$^2$.}}
\gap
\begin{probbox}{48}{\stars{2}}
A square has a diagonal of $8$ cm. Find its area.
\end{probbox}
\probsol{
A diagonal splits the square into two equal triangles on the base $8$ cm.\par The diagonals are equal, perpendicular and bisect each other, so each triangle's height is $8:2=4$ cm.\par Each triangle: $8\times4:2=16$ cm$^2$.\par Square: $32$ cm$^2$.
\dapso{$32$ cm$^2$.}}
\gap
\begin{probbox}{49}{\stars{2}}
Triangle ABC has area $80$ cm$^2$. M is the midpoint of BC and N the midpoint of AM. Find the area of BNC.
\end{probbox}
\probsol{
$S_{ABM}=S_{ACM}=40$ cm$^2$.\par BNM and BAM share the height from B and $NM=\dfrac12AM$, so $S_{BNM}=20$ cm$^2$.\par Likewise $S_{CNM}=20$ cm$^2$.\par $S_{BNC}=40$ cm$^2$.
\dapso{$40$ cm$^2$.}}
\gap
\begin{probbox}{50}{\stars{2}}
Triangle ABC has area $120$ cm$^2$. M is on BC with $BM=\dfrac14BC$; N is on AM with $AN=\dfrac23AM$. Find the areas of ABN and NBM.
\end{probbox}
\probsol{
$S_{ABM}=120:4=30$ cm$^2$ (same height from A).\par ABN and ABM share the height from B, with $AN=\dfrac23AM$.\par $S_{ABN}=20$ cm$^2$.\par $S_{NBM}=10$ cm$^2$.
\dapso{$S_{ABN}=20$ cm$^2$, $S_{NBM}=10$ cm$^2$.}}
\clearpage
\titlehead
\begingroup\setbox0=\vbox{\section*{(*) Comparing Areas of Triangles}}\endgroup
\chapterhead{12}{(*) Comparing Areas of Triangles}

\shead{Same height, compare the bases}
\begin{knowbox}
\begin{minipage}[t]{0.6\linewidth}
\begin{itemize}[leftmargin=1.4em,itemsep=3pt]
\item Two triangles with \textbf{the same height} have areas in the same ratio as \textbf{their bases}. For example, if M is on BC then
\[\frac{S_{ABM}}{S_{AMC}}=\frac{BM}{MC}.\]
\item Two triangles with \textbf{the same base} have areas in the same ratio as \textbf{their heights}.
\item The segment from a vertex to the \textbf{midpoint} of the opposite side splits a triangle into two triangles of equal area.
\end{itemize}
To compare two triangles, it often helps to pass through a ``middle'' triangle that shares a height with each of them.
\end{minipage}\hfill
\begin{minipage}[t]{0.36\linewidth}\vspace{4pt}\centering
\begin{tikzpicture}[scale=0.85]
\coordinate (B) at (0,0); \coordinate (C) at (4.2,0); \coordinate (A) at (1.4,2.6); \coordinate (M) at (2.8,0); \coordinate (H) at (1.4,0);
\fill[tmBlue!18] (A)--(B)--(M)--cycle; \fill[tmOrange!18] (A)--(M)--(C)--cycle;
\draw[geo] (A)--(B)--(C)--cycle; \draw[geo] (A)--(M);
\draw[dashed,Slate] (A)--(H); \ramark{H}{C}{A}
\node[lbl,above] at (A){A}; \node[lbl,below] at (B){B}; \node[lbl,below] at (C){C}; \node[lbl,below] at (M){M};
\end{tikzpicture}
\end{minipage}
\end{knowbox}
\gap
\begin{exbox}{\ExWord\ 1}
Triangle ABC has area $45$ cm$^2$. Point M is on BC with $BM$ twice $MC$. Find the areas of triangles ABM and AMC.
\guide
Triangles ABM and AMC share the height from A, so their areas are in the ratio $BM:MC=2:1$.\par If $S_{AMC}$ is $1$ part, $S_{ABM}$ is $2$ parts and $S_{ABC}$ is $3$ parts.\par
$S_{AMC}=45:3=15$ cm$^2$;\par $S_{ABM}=15\times2=30$ cm$^2$.
\dapso{$30$ cm$^2$ and $15$ cm$^2$.}
\end{exbox}
\gap
\begin{exbox}{\ExWord\ 2}
\begin{minipage}[t]{0.6\linewidth}
Triangle ABC has area $60$ cm$^2$; M is the midpoint of BC and N is the midpoint of AM. Find the areas of triangles ABN and BNC.
\end{minipage}\hfill
\begin{minipage}[t]{0.36\linewidth}\vspace{0pt}\centering
\begin{tikzpicture}[scale=0.7]
\coordinate (B) at (0,0); \coordinate (C) at (4,0); \coordinate (A) at (1.2,2.6); \coordinate (M) at (2,0); \coordinate (N) at (1.6,1.3);
\fill[shade] (B)--(N)--(C)--cycle;
\draw[geo] (A)--(B)--(C)--cycle; \draw[geo] (A)--(M) (B)--(N)--(C);
\node[lbl,above] at (A){A}; \node[lbl,below] at (B){B}; \node[lbl,below] at (C){C}; \node[lbl,below] at (M){M}; \node[lbl,right] at (N){N};
\end{tikzpicture}
\end{minipage}
\guide
M is the midpoint of BC, so $S_{ABM}=S_{AMC}=60:2=30$ cm$^2$.\par
N is the midpoint of AM, so $S_{ABN}=S_{NBM}=30:2=15$ cm$^2$; likewise $S_{NMC}=15$ cm$^2$.\par
$S_{BNC}=S_{NBM}+S_{NMC}=15+15=30$ cm$^2$.
\dapso{$S_{ABN}=15$ cm$^2$;\par $S_{BNC}=30$ cm$^2$.}
\end{exbox}
\gap
\begin{exbox}{\ExWord\ 3}
Triangle ABC has area $72$ cm$^2$. M is on AB with $AM=\dfrac13AB$, and N is on AC with $AN=\dfrac12AC$. Find the areas of triangle AMN and quadrilateral BMNC.
\guide
Join M to C (the ``middle'' triangle AMC).\par
AMC and ABC share the height from C, and $AM=\dfrac13AB$, so $S_{AMC}=\dfrac13\times72=24$ cm$^2$.\par
AMN and AMC share the height from M, and $AN=\dfrac12AC$, so $S_{AMN}=\dfrac12\times24=12$ cm$^2$.\par
$S_{BMNC}=72-12=60$ cm$^2$.
\dapso{$12$ cm$^2$; $60$ cm$^2$.}
\end{exbox}
\gap
\begin{exbox}{\ExWord\ 4}
Square ABCD has side $12$ cm; M is the midpoint of AB and N the midpoint of BC. Find the area of triangle DMN.
\guide
Take the square and remove the three right-angled corner triangles:\par
$S_{AMD}=6\times12:2=36$;\par $S_{MBN}=6\times6:2=18$;\par $S_{NCD}=6\times12:2=36$ cm$^2$.\par
$S_{DMN}=12\times12-36-18-36=54$ cm$^2$.
\dapso{$54$ cm$^2$.}
\end{exbox}
\dividerbar{Practice}
\begin{probbox}{1}{\relax}
Triangle ABC has area $40$ cm$^2$. M is on BC with $BM=3$ cm and $MC=5$ cm. Find the areas of ABM and AMC.
\end{probbox}
\probsol{
The areas are in the ratio $3:5$, $8$ parts in all.\par $S_{ABM}=40:8\times3=15$ cm$^2$; $S_{AMC}=25$ cm$^2$.
\dapso{$15$ cm$^2$ and $25$ cm$^2$.}}
\gap
\begin{probbox}{2}{\relax}
In triangle ABC, M is the midpoint of BC and $S_{ABM}=18$ cm$^2$. Find $S_{ABC}$.
\end{probbox}
\probsol{
$S_{AMC}=S_{ABM}=18$ cm$^2$, so $S_{ABC}=36$ cm$^2$.
\dapso{$36$ cm$^2$.}}
\gap
\begin{probbox}{3}{\relax}
Triangle ABC has area $90$ cm$^2$. D is on AC with $AD=\dfrac23AC$. Find the areas of ABD and BDC.
\end{probbox}
\probsol{
They share the height from B: $S_{ABD}=\dfrac23\times90=60$ cm$^2$ and $S_{BDC}=30$ cm$^2$.
\dapso{$60$ cm$^2$ and $30$ cm$^2$.}}
\gap
\begin{probbox}{4}{\relax}
Points M and N divide side BC of triangle ABC into three equal parts $BM=MN=NC$. If $S_{AMN}=15$ cm$^2$, find $S_{ABC}$.
\end{probbox}
\probsol{
ABM, AMN and ANC share a height and have equal bases, so their areas are equal.\par $S_{ABC}=15\times3=45$ cm$^2$.
\dapso{$45$ cm$^2$.}}
\gap
\begin{probbox}{5}{\relax}
Triangle ABC has area $20$ cm$^2$. Extend BC beyond C to D with $CD=BC$. Find the areas of ACD and ABD.
\end{probbox}
\probsol{
ACD and ABC share the height from A and have equal bases: $S_{ACD}=20$ cm$^2$.\par $S_{ABD}=20+20=40$ cm$^2$.
\dapso{$20$ cm$^2$; $40$ cm$^2$.}}
\gap
\begin{probbox}{6}{\relax}
Triangle ABC has area $60$ cm$^2$; M is the midpoint of BC; N is on AC with $AN=\dfrac13AC$. Find $S_{MNC}$.
\end{probbox}
\probsol{
$S_{AMC}=60:2=30$ cm$^2$.\par MNC and AMC share the height from M and $NC=\dfrac23AC$, so $S_{MNC}=\dfrac23\times30=20$ cm$^2$.
\dapso{$20$ cm$^2$.}}
\gap
\begin{probbox}{7}{\relax}
Triangle ABC has area $54$ cm$^2$. M is on BC with $BM=\dfrac13BC$; N is the midpoint of AM. Find the areas of BNM and ANC.
\end{probbox}
\probsol{
$S_{ABM}=\dfrac13\times54=18$ and $S_{AMC}=36$ cm$^2$.\par N is the midpoint of AM, so $S_{BNM}=18:2=9$ cm$^2$ and $S_{ANC}=36:2=18$ cm$^2$.
\dapso{$9$ cm$^2$; $18$ cm$^2$.}}
\gap
\begin{probbox}{8}{\relax}
Triangle ABC has area $48$ cm$^2$. M is on AB with $AM=\dfrac23AB$; N is on AC with $AN=\dfrac34AC$. Find the areas of AMN and BMNC.
\end{probbox}
\probsol{
$S_{AMC}=\dfrac23\times48=32$ cm$^2$;\par $S_{AMN}=\dfrac34\times32=24$ cm$^2$;\par $S_{BMNC}=48-24=24$ cm$^2$.
\dapso{$24$ cm$^2$; $24$ cm$^2$.}}
\gap
\begin{probbox}{9}{\relax}
In triangle ABC, D is on BC with $BD=\dfrac14BC$, and $S_{ABD}=12$ cm$^2$. Find $S_{ADC}$.
\end{probbox}
\probsol{
$DC=3\times BD$ and they share the height from A, so $S_{ADC}=12\times3=36$ cm$^2$.
\dapso{$36$ cm$^2$.}}
\gap
\begin{probbox}{10}{\relax}
Triangle ABC has area $60$ cm$^2$. M on BC divides it so that $BM:MC=1:2$. Find the areas of ABM and AMC.
\end{probbox}
\probsol{
ABM, AMC and ABC share the height from A.\par $BM=\dfrac{1}{3}BC$, so $S_{ABM}=60\times\dfrac{1}{3}=20$ cm$^2$.\par $S_{AMC}=60-20=40$ cm$^2$.
\dapso{$S_{ABM}=20$ cm$^2$, $S_{AMC}=40$ cm$^2$.}}
\gap
\begin{probbox}{11}{\relax}
Triangle ABC has area $84$ cm$^2$. M on BC divides it so that $BM:MC=3:4$. Find the areas of ABM and AMC.
\end{probbox}
\probsol{
ABM, AMC and ABC share the height from A.\par $BM=\dfrac{3}{7}BC$, so $S_{ABM}=84\times\dfrac{3}{7}=36$ cm$^2$.\par $S_{AMC}=84-36=48$ cm$^2$.
\dapso{$S_{ABM}=36$ cm$^2$, $S_{AMC}=48$ cm$^2$.}}
\gap
\begin{probbox}{12}{\relax}
Triangle ABC has area $100$ cm$^2$. M on BC divides it so that $BM:MC=2:3$. Find the areas of ABM and AMC.
\end{probbox}
\probsol{
ABM, AMC and ABC share the height from A.\par $BM=\dfrac{2}{5}BC$, so $S_{ABM}=100\times\dfrac{2}{5}=40$ cm$^2$.\par $S_{AMC}=100-40=60$ cm$^2$.
\dapso{$S_{ABM}=40$ cm$^2$, $S_{AMC}=60$ cm$^2$.}}
\gap
\begin{probbox}{13}{\relax}
Triangle ABC has area $72$ cm$^2$. M on BC divides it so that $BM:MC=5:1$. Find the areas of ABM and AMC.
\end{probbox}
\probsol{
ABM, AMC and ABC share the height from A.\par $BM=\dfrac{5}{6}BC$, so $S_{ABM}=72\times\dfrac{5}{6}=60$ cm$^2$.\par $S_{AMC}=72-60=12$ cm$^2$.
\dapso{$S_{ABM}=60$ cm$^2$, $S_{AMC}=12$ cm$^2$.}}
\gap
\begin{probbox}{14}{\relax}
Triangle ABC has area $45$ cm$^2$. M on BC divides it so that $BM:MC=4:5$. Find the areas of ABM and AMC.
\end{probbox}
\probsol{
ABM, AMC and ABC share the height from A.\par $BM=\dfrac{4}{9}BC$, so $S_{ABM}=45\times\dfrac{4}{9}=20$ cm$^2$.\par $S_{AMC}=45-20=25$ cm$^2$.
\dapso{$S_{ABM}=20$ cm$^2$, $S_{AMC}=25$ cm$^2$.}}
\gap
\begin{probbox}{15}{\relax}
In triangle ABC, M on BC has $BM=\dfrac{1}{3}BC$. Triangle ABM has area $15$ cm$^2$. Find the areas of ABC and AMC.
\end{probbox}
\probsol{
ABM and ABC share the height from A, so $S_{ABM}=\dfrac{1}{3}S_{ABC}$.\par $S_{ABC}=15:\dfrac{1}{3}=45$ cm$^2$.\par $S_{AMC}=45-15=30$ cm$^2$.
\dapso{$S_{ABC}=45$ cm$^2$, $S_{AMC}=30$ cm$^2$.}}
\gap
\begin{probbox}{16}{\relax}
In triangle ABC, M on BC has $BM=\dfrac{2}{5}BC$. Triangle ABM has area $24$ cm$^2$. Find the areas of ABC and AMC.
\end{probbox}
\probsol{
ABM and ABC share the height from A, so $S_{ABM}=\dfrac{2}{5}S_{ABC}$.\par $S_{ABC}=24:\dfrac{2}{5}=60$ cm$^2$.\par $S_{AMC}=60-24=36$ cm$^2$.
\dapso{$S_{ABC}=60$ cm$^2$, $S_{AMC}=36$ cm$^2$.}}
\gap
\begin{probbox}{17}{\relax}
In triangle ABC, M on BC has $BM=\dfrac{3}{4}BC$. Triangle ABM has area $21$ cm$^2$. Find the areas of ABC and AMC.
\end{probbox}
\probsol{
ABM and ABC share the height from A, so $S_{ABM}=\dfrac{3}{4}S_{ABC}$.\par $S_{ABC}=21:\dfrac{3}{4}=28$ cm$^2$.\par $S_{AMC}=28-21=7$ cm$^2$.
\dapso{$S_{ABC}=28$ cm$^2$, $S_{AMC}=7$ cm$^2$.}}
\gap
\begin{probbox}{18}{\relax}
In triangle ABC, M on BC has $BM=\dfrac{1}{6}BC$. Triangle ABM has area $12$ cm$^2$. Find the areas of ABC and AMC.
\end{probbox}
\probsol{
ABM and ABC share the height from A, so $S_{ABM}=\dfrac{1}{6}S_{ABC}$.\par $S_{ABC}=12:\dfrac{1}{6}=72$ cm$^2$.\par $S_{AMC}=72-12=60$ cm$^2$.
\dapso{$S_{ABC}=72$ cm$^2$, $S_{AMC}=60$ cm$^2$.}}
\gap
\begin{probbox}{19}{\relax}
Triangle ABC has area $80$ cm$^2$ and $BC=16$ cm. Point M on BC makes triangle ABM have area $30$ cm$^2$. Find BM.
\end{probbox}
\probsol{
ABM and ABC share the height from A, so their areas are in the ratio of the bases.\par $\dfrac{BM}{BC}=\dfrac{30}{80}=\dfrac{3}{8}$.\par $BM=16\times\dfrac{3}{8}=6$ cm.
\dapso{$BM=6$ cm.}}
\gap
\begin{probbox}{20}{\relax}
Triangle ABC has area $96$ cm$^2$ and $BC=12$ cm. Point M on BC makes triangle ABM have area $36$ cm$^2$. Find BM.
\end{probbox}
\probsol{
ABM and ABC share the height from A, so their areas are in the ratio of the bases.\par $\dfrac{BM}{BC}=\dfrac{36}{96}=\dfrac{3}{8}$.\par $BM=12\times\dfrac{3}{8}=4.5$ cm.
\dapso{$BM=4.5$ cm.}}
\gap
\begin{probbox}{21}{\relax}
Triangle ABC has area $75$ cm$^2$ and $BC=10$ cm. Point M on BC makes triangle ABM have area $45$ cm$^2$. Find BM.
\end{probbox}
\probsol{
ABM and ABC share the height from A, so their areas are in the ratio of the bases.\par $\dfrac{BM}{BC}=\dfrac{45}{75}=\dfrac{3}{5}$.\par $BM=10\times\dfrac{3}{5}=6$ cm.
\dapso{$BM=6$ cm.}}
\gap
\begin{probbox}{22}{\relax}
Triangle ABC has area $36$ cm$^2$. Extend BC beyond C by $CD=\dfrac12BC$. Find the areas of ACD and ABD.
\end{probbox}
\probsol{
ABC, ACD and ABD share the height from A.\par $S_{ACD}=36\times\dfrac{1}{2}=18$ cm$^2$.\par $S_{ABD}=36+18=54$ cm$^2$.
\dapso{$S_{ACD}=18$ cm$^2$, $S_{ABD}=54$ cm$^2$.}}
\gap
\begin{probbox}{23}{\relax}
Triangle ABC has area $50$ cm$^2$. Extend BC beyond C by $CD=\dfrac25BC$. Find the areas of ACD and ABD.
\end{probbox}
\probsol{
ABC, ACD and ABD share the height from A.\par $S_{ACD}=50\times\dfrac{2}{5}=20$ cm$^2$.\par $S_{ABD}=50+20=70$ cm$^2$.
\dapso{$S_{ACD}=20$ cm$^2$, $S_{ABD}=70$ cm$^2$.}}
\gap
\begin{probbox}{24}{\relax}
Triangle ABC has area $28$ cm$^2$. Extend BC beyond C by $CD=2\times BC$. Find the areas of ACD and ABD.
\end{probbox}
\probsol{
ABC, ACD and ABD share the height from A.\par $S_{ACD}=28\times2=56$ cm$^2$.\par $S_{ABD}=28+56=84$ cm$^2$.
\dapso{$S_{ACD}=56$ cm$^2$, $S_{ABD}=84$ cm$^2$.}}
\gap
\begin{probbox}{25}{\relax}
Triangle ABC has area $60$ cm$^2$. M on AB has $AM=\dfrac{1}{2}AB$; N on AC has $AN=\dfrac{1}{3}AC$. Find the areas of AMN and of quadrilateral BMNC.
\end{probbox}
\probsol{
Join M to C.\par AMC and ABC share the height from C: $S_{AMC}=60\times\dfrac{1}{2}=30$ cm$^2$.\par AMN and AMC share the height from M: $S_{AMN}=30\times\dfrac{1}{3}=10$ cm$^2$.\par $S_{BMNC}=60-10=50$ cm$^2$.
\dapso{$S_{AMN}=10$ cm$^2$, $S_{BMNC}=50$ cm$^2$.}}
\gap
\begin{probbox}{26}{\relax}
Triangle ABC has area $90$ cm$^2$. M on AB has $AM=\dfrac{2}{3}AB$; N on AC has $AN=\dfrac{1}{2}AC$. Find the areas of AMN and of quadrilateral BMNC.
\end{probbox}
\probsol{
Join M to C.\par AMC and ABC share the height from C: $S_{AMC}=90\times\dfrac{2}{3}=60$ cm$^2$.\par AMN and AMC share the height from M: $S_{AMN}=60\times\dfrac{1}{2}=30$ cm$^2$.\par $S_{BMNC}=90-30=60$ cm$^2$.
\dapso{$S_{AMN}=30$ cm$^2$, $S_{BMNC}=60$ cm$^2$.}}
\gap
\begin{probbox}{27}{\relax}
Triangle ABC has area $84$ cm$^2$. M on AB has $AM=\dfrac{3}{4}AB$; N on AC has $AN=\dfrac{2}{3}AC$. Find the areas of AMN and of quadrilateral BMNC.
\end{probbox}
\probsol{
Join M to C.\par AMC and ABC share the height from C: $S_{AMC}=84\times\dfrac{3}{4}=63$ cm$^2$.\par AMN and AMC share the height from M: $S_{AMN}=63\times\dfrac{2}{3}=42$ cm$^2$.\par $S_{BMNC}=84-42=42$ cm$^2$.
\dapso{$S_{AMN}=42$ cm$^2$, $S_{BMNC}=42$ cm$^2$.}}
\gap
\begin{probbox}{28}{\relax}
Triangles ABC and DBC share the base BC. The height from A is $8$ cm and from D is $6$ cm. If $S_{ABC}=40$ cm$^2$, find $S_{DBC}$.
\end{probbox}
\probsol{
With a shared base, areas are in the ratio of the heights: $\dfrac{S_{DBC}}{S_{ABC}}=\dfrac{6}{8}$.\par $S_{DBC}=40\times6:8=30$ cm$^2$.
\dapso{$30$ cm$^2$.}}
\gap
\begin{probbox}{29}{\relax}
Triangles ABC and DBC share the base BC. The height from A is $9$ cm and from D is $12$ cm. If $S_{ABC}=54$ cm$^2$, find $S_{DBC}$.
\end{probbox}
\probsol{
With a shared base, areas are in the ratio of the heights: $\dfrac{S_{DBC}}{S_{ABC}}=\dfrac{12}{9}$.\par $S_{DBC}=54\times12:9=72$ cm$^2$.
\dapso{$72$ cm$^2$.}}
\gap
\begin{probbox}{30}{\relax}
Triangles ABC and DBC share the base BC. The height from A is $10$ cm and from D is $4$ cm. If $S_{ABC}=35$ cm$^2$, find $S_{DBC}$.
\end{probbox}
\probsol{
With a shared base, areas are in the ratio of the heights: $\dfrac{S_{DBC}}{S_{ABC}}=\dfrac{4}{10}$.\par $S_{DBC}=35\times4:10=14$ cm$^2$.
\dapso{$14$ cm$^2$.}}
\dividerbar{Challenge -- Extension}
\begin{probbox}{31}{\stars{1}}
Triangle ABC has area $64$ cm$^2$; M and N are the midpoints of AB and AC. Find the areas of AMN and BMNC.
\end{probbox}
\probsol{
$S_{AMC}=64:2=32$ cm$^2$; $S_{AMN}=32:2=16$ cm$^2$.\par $S_{BMNC}=64-16=48$ cm$^2$.
\note{The area of AMN is $\dfrac14$ of the area of ABC.}
\dapso{$16$ cm$^2$; $48$ cm$^2$.}}
\gap
\begin{probbox}{32}{\stars{1}}
Triangle ABC has area $30$ cm$^2$. Extend AB beyond B to D with $BD=AB$, and AC beyond C to E with $CE=AC$. Find $S_{ADE}$.
\end{probbox}
\probsol{
$AD=2\times AB$, so $S_{ADC}=2\times30=60$ cm$^2$.\par $AE=2\times AC$, so $S_{ADE}=2\times60=120$ cm$^2$.
\dapso{$120$ cm$^2$.}}
\gap
\begin{probbox}{33}{\stars{1}}
Triangle ABC has area $120$ cm$^2$. D is on BC with $BD=2\times DC$; E is the midpoint of AD. Find $S_{BEC}$.
\end{probbox}
\probsol{
$S_{ABD}=80$ and $S_{ADC}=40$ cm$^2$.\par E is the midpoint of AD, so $S_{BED}=40$ and $S_{EDC}=20$ cm$^2$.\par
$S_{BEC}=40+20=60$ cm$^2$.
\dapso{$60$ cm$^2$.}}
\gap
\begin{probbox}{34}{\stars{2}}
\begin{minipage}[t]{0.6\linewidth}
In triangle ABC, M is the midpoint of BC and I is the midpoint of AM. Line BI meets AC at N. Show that $AN=\dfrac13AC$.
\end{minipage}\hfill
\begin{minipage}[t]{0.36\linewidth}\vspace{0pt}\centering
\begin{tikzpicture}[scale=0.7]
\coordinate (B) at (0,0); \coordinate (C) at (4.2,0); \coordinate (A) at (1.2,2.7); \coordinate (M) at (2.1,0); \coordinate (I) at (1.65,1.35); \coordinate (N) at (2.2,1.8);
\draw[geo] (A)--(B)--(C)--cycle; \draw[geo] (A)--(M) (B)--(N);
\draw[dashed,Slate] (M)--(N);
\node[lbl,above] at (A){A}; \node[lbl,below] at (B){B}; \node[lbl,below] at (C){C}; \node[lbl,below] at (M){M}; \node[lbl,left] at (I){I}; \node[lbl,above right] at (N){N};
\end{tikzpicture}
\end{minipage}
\end{probbox}
\probsol{
Join M to N.\par Since $AI=IM$, $S_{ABI}=S_{MBI}$ and $S_{AIN}=S_{MIN}$.\par Adding: $S_{ABN}=S_{MBN}$.\par
Since $BM=MC$, $S_{MBN}=S_{MNC}$.\par Hence $S_{BNC}=2\times S_{ABN}$.\par
ABN and CBN share the height from B, so $NC=2\times AN$, i.e.\par $AN=\dfrac13AC$.
\dapso{Shown as required.}}
\gap
\begin{probbox}{35}{\stars{2}}
Triangle ABC has area $90$ cm$^2$. M is on BC with $BM=\dfrac13BC$, N on CA with $CN=\dfrac13CA$, and P on AB with $AP=\dfrac13AB$. Find $S_{MNP}$.
\end{probbox}
\probsol{
Triangle APN: $S_{APC}=\dfrac13\times90=30$; $AN=\dfrac23AC$, so $S_{APN}=\dfrac23\times30=20$ cm$^2$.\par
Likewise $S_{BMP}=20$ and $S_{CNM}=20$ cm$^2$.\par So $S_{MNP}=90-3\times20=30$ cm$^2$.
\dapso{$30$ cm$^2$.}}
\gap
\begin{probbox}{36}{\stars{1}}
Triangle ABC has area $60$ cm$^2$. M on BC has $BM=\dfrac{1}{3}BC$; N on AM has $AN=\dfrac{1}{2}AM$. Find the areas of ABN, NBM and ANC.
\end{probbox}
\probsol{
$S_{ABM}=60\times\dfrac{1}{3}=20$ cm$^2$; $S_{AMC}=40$ cm$^2$.\par ABN and ABM share the height from B: $S_{ABN}=20\times\dfrac{1}{2}=10$ cm$^2$.\par $S_{NBM}=20-10=10$ cm$^2$.\par ANC and AMC share the height from C: $S_{ANC}=40\times\dfrac{1}{2}=20$ cm$^2$.
\dapso{$S_{ABN}=10$, $S_{NBM}=10$, $S_{ANC}=20$ (cm$^2$).}}
\gap
\begin{probbox}{37}{\stars{1}}
Triangle ABC has area $96$ cm$^2$. M on BC has $BM=\dfrac{1}{4}BC$; N on AM has $AN=\dfrac{2}{3}AM$. Find the areas of ABN, NBM and ANC.
\end{probbox}
\probsol{
$S_{ABM}=96\times\dfrac{1}{4}=24$ cm$^2$; $S_{AMC}=72$ cm$^2$.\par ABN and ABM share the height from B: $S_{ABN}=24\times\dfrac{2}{3}=16$ cm$^2$.\par $S_{NBM}=24-16=8$ cm$^2$.\par ANC and AMC share the height from C: $S_{ANC}=72\times\dfrac{2}{3}=48$ cm$^2$.
\dapso{$S_{ABN}=16$, $S_{NBM}=8$, $S_{ANC}=48$ (cm$^2$).}}
\gap
\begin{probbox}{38}{\stars{1}}
Triangle ABC has area $150$ cm$^2$. M on BC has $BM=\dfrac{2}{5}BC$; N on AM has $AN=\dfrac{3}{5}AM$. Find the areas of ABN, NBM and ANC.
\end{probbox}
\probsol{
$S_{ABM}=150\times\dfrac{2}{5}=60$ cm$^2$; $S_{AMC}=90$ cm$^2$.\par ABN and ABM share the height from B: $S_{ABN}=60\times\dfrac{3}{5}=36$ cm$^2$.\par $S_{NBM}=60-36=24$ cm$^2$.\par ANC and AMC share the height from C: $S_{ANC}=90\times\dfrac{3}{5}=54$ cm$^2$.
\dapso{$S_{ABN}=36$, $S_{NBM}=24$, $S_{ANC}=54$ (cm$^2$).}}
\gap
\begin{probbox}{39}{\stars{1}}
M on AB has $AM=\dfrac{1}{2}AB$ and N on AC has $AN=\dfrac{2}{5}AC$. Triangle AMN has area $10$ cm$^2$. Find the area of ABC.
\end{probbox}
\probsol{
Join M to C.\par AMN and AMC share the height from M: $S_{AMC}=10:\dfrac{2}{5}=25$ cm$^2$.\par AMC and ABC share the height from C: $S_{ABC}=25:\dfrac{1}{2}=50$ cm$^2$.
\dapso{$50$ cm$^2$.}}
\gap
\begin{probbox}{40}{\stars{1}}
M on AB has $AM=\dfrac{3}{4}AB$ and N on AC has $AN=\dfrac{2}{3}AC$. Triangle AMN has area $18$ cm$^2$. Find the area of ABC.
\end{probbox}
\probsol{
Join M to C.\par AMN and AMC share the height from M: $S_{AMC}=18:\dfrac{2}{3}=27$ cm$^2$.\par AMC and ABC share the height from C: $S_{ABC}=27:\dfrac{3}{4}=36$ cm$^2$.
\dapso{$36$ cm$^2$.}}
\gap
\begin{probbox}{41}{\stars{1}}
Triangle ABC has area $12$ cm$^2$. Extend AB beyond B by $BD=AB$ and AC beyond C by $CE=2\times AC$. Find the area of ADE.
\end{probbox}
\probsol{
$AD=2\times AB$, $AE=3\times AC$.\par Join D to C.\par ADC and ABC share the height from C: $S_{ADC}=12\times2=24$ cm$^2$.\par ADE and ADC share the height from D: $S_{ADE}=24\times3=72$ cm$^2$.
\dapso{$72$ cm$^2$.}}
\gap
\begin{probbox}{42}{\stars{1}}
Triangle ABC has area $15$ cm$^2$. Extend AB beyond B by $BD=2\times AB$ and AC beyond C by $CE=2\times AC$. Find the area of ADE.
\end{probbox}
\probsol{
$AD=3\times AB$, $AE=3\times AC$.\par Join D to C.\par ADC and ABC share the height from C: $S_{ADC}=15\times3=45$ cm$^2$.\par ADE and ADC share the height from D: $S_{ADE}=45\times3=135$ cm$^2$.
\dapso{$135$ cm$^2$.}}
\gap
\begin{probbox}{43}{\stars{2}}
Triangle ABC has area $80$ cm$^2$. M, N, P lie on AB, BC, CA with $AM=\dfrac{1}{4}AB$, $BN=\dfrac{1}{4}BC$, $CP=\dfrac{1}{4}CA$. Find the area of MNP.
\end{probbox}
\probsol{
In triangle AMP: $AM=\dfrac{1}{4}AB$, $AP=\dfrac{3}{4}AC$.\par As before (join M to C): $S_{AMP}=80\times\dfrac{1}{4}\times\dfrac{3}{4}=15$ cm$^2$.\par Likewise $S_{BNM}=S_{CPN}=15$ cm$^2$.\par $S_{MNP}=80-3\times15=35$ cm$^2$.
\dapso{$35$ cm$^2$.}}
\gap
\begin{probbox}{44}{\stars{2}}
Triangle ABC has area $72$ cm$^2$. M, N, P lie on AB, BC, CA with $AM=\dfrac{1}{3}AB$, $BN=\dfrac{1}{3}BC$, $CP=\dfrac{1}{3}CA$. Find the area of MNP.
\end{probbox}
\probsol{
In triangle AMP: $AM=\dfrac{1}{3}AB$, $AP=\dfrac{2}{3}AC$.\par As before (join M to C): $S_{AMP}=72\times\dfrac{1}{3}\times\dfrac{2}{3}=16$ cm$^2$.\par Likewise $S_{BNM}=S_{CPN}=16$ cm$^2$.\par $S_{MNP}=72-3\times16=24$ cm$^2$.
\dapso{$24$ cm$^2$.}}
\gap
\begin{probbox}{45}{\stars{1}}
In triangle ABC, M on BC has $BM:MC=3:1$. Triangle ABM is $20$ cm$^2$ larger than AMC. Find the area of ABC.
\end{probbox}
\probsol{
ABM and AMC share the height from A, so $S_{ABM}:S_{AMC}=3:1$.\par The difference of $2$ parts is $20$ cm$^2$; one part is $10$ cm$^2$.\par $S_{ABC}=10\times4=40$ cm$^2$.
\dapso{$40$ cm$^2$.}}
\gap
\begin{probbox}{46}{\stars{1}}
In triangle ABC, M on BC has $BM:MC=5:2$. Triangle ABM is $24$ cm$^2$ larger than AMC. Find the area of ABC.
\end{probbox}
\probsol{
ABM and AMC share the height from A, so $S_{ABM}:S_{AMC}=5:2$.\par The difference of $3$ parts is $24$ cm$^2$; one part is $8$ cm$^2$.\par $S_{ABC}=8\times7=56$ cm$^2$.
\dapso{$56$ cm$^2$.}}
\gap
\begin{probbox}{47}{\stars{1}}
Rectangle ABCD has $AB=24$ cm, $BC=16$ cm. M on AB has $AM=8$ cm; N on BC has $BN=6$ cm. Find the area of triangle DMN.
\end{probbox}
\probsol{
$S_{ABCD}=24\times16=384$ cm$^2$.\par $S_{AMD}=64$, $S_{MBN}=48$, $S_{NCD}=120$ cm$^2$.\par $S_{DMN}=384-64-48-120=152$ cm$^2$.
\dapso{$152$ cm$^2$.}}
\gap
\begin{probbox}{48}{\stars{1}}
Rectangle ABCD has $AB=30$ cm, $BC=20$ cm. M on AB has $AM=12$ cm; N on BC has $BN=10$ cm. Find the area of triangle DMN.
\end{probbox}
\probsol{
$S_{ABCD}=30\times20=600$ cm$^2$.\par $S_{AMD}=120$, $S_{MBN}=90$, $S_{NCD}=150$ cm$^2$.\par $S_{DMN}=600-120-90-150=240$ cm$^2$.
\dapso{$240$ cm$^2$.}}
\gap
\begin{probbox}{49}{\stars{1}}
Rectangle ABCD has $AB=18$ cm, $BC=12$ cm. M on AB has $AM=6$ cm; N on BC has $BN=4$ cm. Find the area of triangle DMN.
\end{probbox}
\probsol{
$S_{ABCD}=18\times12=216$ cm$^2$.\par $S_{AMD}=36$, $S_{MBN}=24$, $S_{NCD}=72$ cm$^2$.\par $S_{DMN}=216-36-24-72=84$ cm$^2$.
\dapso{$84$ cm$^2$.}}
\gap
\begin{probbox}{50}{\stars{1}}
Triangle ABC has $BC=16$ cm. Extending BC beyond C by $5$ cm gives a triangle $25$ cm$^2$ larger. Find the area of ABC.
\end{probbox}
\probsol{
The extra part is a triangle with base $5$ cm and the same height from A.\par Height: $25\times2:5=10$ cm.\par $S_{ABC}=16\times10:2=80$ cm$^2$.
\dapso{$80$ cm$^2$.}}
\clearpage
\titlehead
